% ============ Appendix A. Tensor products, direct and inverse limits (原书 pp.266-273) ============
\chapter*{Appendix A}
\addcontentsline{toc}{chapter}{Appendix A: Tensor products, direct and inverse limits}
\markboth{Appendix A}{Tensor products, direct and inverse limits}
\renewcommand{\thetheorem}{A\arabic{theorem}}
\setcounter{theorem}{0}

\vspace*{-0.5cm}
\begin{center}
  {\Large\bfseries Tensor products, direct and inverse limits}
\end{center}
\vspace{0.3cm}

\medskip
\noindent\textbf{Tensor products}

Let $A$ be a ring, $L$, $M$ and $N$ three $A$-modules. We say that a map
$\varphi\colon M \times N \to L$ is \emph{bilinear} if fixing either of the entries it is $A$-linear
in the other, that is if
\[
  \varphi(x + x', y) = \varphi(x, y) + \varphi(x', y), \quad \varphi(ax, y) = a\varphi(x, y),
\]
\[
  \varphi(x, y + y') = \varphi(x, y) + \varphi(x, y'), \quad \varphi(x, ay) = a\varphi(x, y).
\]
Write $\mathcal{B}(M, N; L)$ or $\mathcal{B}_A(M, N; L)$ for the set of all bilinear maps from
$M \times N$ to $L$; as with $\Hom(M, L)$, this has an $A$-module structure (since we are
assuming that $A$ is commutative).

If $g\colon L \to L'$ is an $A$-linear map and $\varphi \in \mathcal{B}(M, N; L)$ then
$g\varphi \in \mathcal{B}(M, N; L')$. Bearing this in mind, for given $M$ and $N$, consider a
bilinear map $\theta\colon M \times N \to L_0$ having the following property, where we write
$x \otimes y$ instead of $\theta(x, y)$: for any $A$-module $L$ and any $\varphi \in \mathcal{B}(M, N; L)$
there exists a unique $A$-linear map $g\colon L_0 \to L$ satisfying
\[
  \varphi(x \otimes y) = g(x \otimes y).
\]
If this holds we say that $L_0$ is the \emph{tensor product} of $M$ and $N$ over $A$, and
write $L_0 = M \otimes_A N$; we sometimes omit $A$ and write $M \otimes N$. As usual with
this kind of definition, $M \otimes_A N$, assuming it exists, is uniquely determined
up to isomorphism. To prove existence, write $F$ for the free $A$-module with
basis the set $M \times N$, and let $R \subset F$ be the submodule generated by all
elements of the form
\[
  (x + x', y) - (x, y) - (x', y), \quad (ax, y) - a(x, y)
\]
\[
  (x, y + y') - (x, y) - (x, y'), \quad (x, ay) - a(x, y).
\]
Then set $L_0 = F/R$ and write $x \otimes y$ for the image in $L_0$ of $(x, y) \in F$. It is now
easy to check that $L_0$ and $\otimes$ satisfy the above condition.

Note that the general element of $M \otimes_A N$ is a sum of the form $\sum x_i \otimes y_i$,
and cannot necessarily be written $x \otimes y$.

For $A$-modules $M$, $N$ and $L$ the definition of tensor product gives:

\textbf{Formula 1.} $\Hom_A(M \otimes_A N, L) \cong \mathcal{B}(M, N; L)$.

The canonical isomorphism is obtained by taking an element $\varphi$ of the
right-hand side to the element $g$ of the left-hand side satisfying
$g(x \otimes y) = \varphi(x, y)$.

We can define multilinear maps from an $r$-fold product of $A$-modules
$M_1, \ldots, M_r$ to an $A$-module $L$ just as in the bilinear case, and get modules
$\mathcal{B}(M_1, \ldots, M_r; L)$ and $M_1 \otimes_A \cdots \otimes_A M_r$; the following `associative
law' then holds:

\textbf{Formula 2.} $(M \otimes_A M') \otimes_A M'' = M \otimes_A (M' \otimes_A M'')$.

For example, for the first equality it is enough to check that the trilinear map
$M \times M' \times M'' \to (M \otimes_A M') \otimes_A M''$ given by
$(x, y, z) \mapsto (x \otimes y) \otimes z$ has the universal property for trilinear maps, and this
is easy. The following Formulas 3, 4 and 5 are also easy:

\textbf{Formula 3.} $M \otimes_A N \cong N \otimes_A M$ \quad (by $x \otimes y \mapsto y \otimes x$).

\textbf{Formula 4.} $M \otimes_A A = M$.

\textbf{Formula 5.} $(\bigoplus_\lambda M_\lambda) \otimes_A N = \bigoplus_\lambda(M_\lambda \otimes_A N)$.

If $f\colon M \to M'$ and $g\colon N \to N'$ are both $A$-linear then
$(x, y) \mapsto f(x) \otimes g(y)$ is a bilinear map from $M \times N$ to $M' \otimes_A N'$, and so it
defines a linear map $M \otimes_A N \to M' \otimes_A N'$, which we denote $f \otimes g$. From
the definition we have:

\textbf{Formula 6.} $(f \otimes g)(\sum x_i \otimes y_i) = \sum f(x_i) \otimes g(y_i)$.

In particular, if both $f$ and $g$ are surjective then we see from this that
$f \otimes g$ is surjective; its kernel is generated by
$\{x \otimes y \mid f(x) = 0 \text{ or } g(y) = 0\}$. Indeed, let $T \subset M \otimes N$ be the submodule
generated by this set; then $T \subset \Ker(f \otimes g)$ so that $f \otimes g$ induces a linear
map $\alpha\colon (M \otimes N)/T \to M' \otimes N'$; furthermore, we can define a bilinear map
$M' \times N' \to (M \otimes N)/T$ by
\[
  (x', y') \mapsto (x \otimes y \bmod T), \quad \text{where } f(x) = x',\ g(y) = y',
\]
since a different choice of inverse images $x$ and $y$ leads to a difference
belonging to $T$. This defines a linear map $\beta\colon M' \otimes N' \to (M \otimes N)/T$,
which is obviously an inverse of $\alpha$. We summarise the above (writing $1$ for
the identity maps):

\textbf{Formula 7.} Suppose given exact sequences
\[
  0 \to K \to M \xrightarrow{f} M' \to 0 \quad \text{and} \quad 0 \to L \xrightarrow{j} N \to N' \to 0,
\]
then $M' \otimes N' \cong (M \otimes N)/T$, where
\[
  T = (i \otimes 1)(K \otimes N) + (1 \otimes j)(M \otimes L).
\]

\textbf{Formula 8 (right-exactness of the tensor product).} If
\[
  M_1 \xrightarrow{f} M_2 \xrightarrow{g} M_3 \to 0
\]
is an exact sequence then so is
\[
  M_1 \otimes_A N \xrightarrow{f \otimes 1} M_2 \otimes_A N \xrightarrow{g \otimes 1} M_3 \otimes_A N \to 0.
\]

In general, even if $f\colon M \to M'$ is injective, $f \otimes 1\colon M \otimes N \to M' \otimes N$ need
not be. (Counter-example: let $A = \mathbb{Z}$ and $N = \mathbb{Z}/n\mathbb{Z}$ for some $n > 1$. Let
$f\colon \mathbb{Z} \to \mathbb{Z}$ be multiplication by $n$; then $\mathbb{Z} \otimes_A N \cong N \ne 0$, but
$f \otimes 1\colon N \to N$ is the zero map, and so is not injective.) However, if $\Im f$ is a
direct summand of $M'$ (in which case we say that the exact sequence
$0 \to M \to M' \to M'/M \to 0$ \emph{splits}), then there is a map $g\colon M' \to M$ such that
$gf = 1$. Then
\[
  (g \otimes 1)(f \otimes 1) = gf \otimes 1 = 1 \otimes 1
\]
is the identity map of $M \otimes N$, and hence $f \otimes 1$ is injective, and the sequence
$0 \to M \otimes N \to M' \otimes N \to (M'/M) \otimes N \to 0$ is split. In particular if $A$ is a
field then any submodule is a direct summand, so that the operation
$\otimes N$ takes exact sequences into exact sequences; in other word, $\otimes N$ is
an exact functor. For an arbitrary ring $A$, an $A$-module $N$ is said to be
\emph{flat} if $\otimes N$ is an exact functor. For more on this see \S\ref{sec:7}.

\medskip
\noindent\textbf{Change of coefficient ring}

Let $A$ and $B$ be rings, and $P$ a two-sided $A$-$B$-module; that is, for $a \in A$,
$b \in B$ and $x \in P$ the products $ax$ and $xb$ are defined, and in addition to the
usual conditions for $A$-modules and $B$-modules we assume that
\[
  (ax)b = a(xb).
\]
Then multiplication by an element $b \in B$ induces an $A$-linear map of $P$ to
itself, which we continue to denote by $b$. This determines a map
$1 \otimes b\colon M \otimes_A P \to M \otimes_A P$ for any $A$-module $M$, and by definition we take
this to be scalar multiplication by $b$ in $M \otimes_A P$; that is, we set
\[
  (\sum y_i \otimes x_i)b = \sum y_i \otimes x_ib \quad \text{for } y_i \in M \text{ and } x_i \in P.
\]
If $N$ is a $B$-module, then for $\varphi \in \Hom_B(P, N)$ we define the product
$\varphi a$ of $\varphi$ and $a \in A$ by
\[
  (\varphi a)(x) = \varphi(ax) \quad \text{for } x \in P;
\]
we have $\varphi a \in \Hom_B(P, N)$, and this makes $\Hom_B(P, N)$ into an $A$-module.

\textbf{Formula 9.} $\Hom_B(M, \Hom_B(P, N)) \cong \Hom_A(M \otimes_A P, N)$.

\textbf{Formula 10.} $(M \otimes_A P) \otimes_B N \cong M \otimes_A (P \otimes_B N)$.

Both of these are easy to prove, and we leave them to the reader.
Formula 10 generalises Formula 2.

Given a ring homomorphism $\lambda\colon A \to B$, we can think of $B$ as a two-sided
$A$-$B$-module by setting $ab = \lambda(a)b$; then for any $A$-module $M$, $M \otimes_A B$ is a
$B$-module, called the \emph{extension of scalars} in $M$ from $A$ to $B$, and written
$M_{(B)}$. For $A$-modules $M$ and $M'$ the following formula holds, so that
tensor product commutes with change of scalars.

\textbf{Formula 11.} $(M \otimes_A B) \otimes_B (M' \otimes_A B) = (M \otimes_A M') \otimes_A B$.

Indeed, using Formulas 10, 4 and 2, the left-hand side is equal to
$M \otimes_A (B \otimes_B (M' \otimes_A B)) = M \otimes_A (M' \otimes_A B) = (M \otimes_A M') \otimes_A B$.

\medskip
\noindent\textbf{Tensor product of $A$-algebras}

Given a ring homomorphism $\lambda\colon A \to B$ we say that $B$ is an \emph{$A$-algebra}. Let
$B'$ be another $A$-algebra defined by $\lambda'\colon A \to B'$. We say that a map
$f\colon B \to B'$ is a homomorphism of $A$-algebras if it is a ring homomorphism
satisfying $\lambda' = f\lambda$. If $B$ and $C$ are $A$-algebras, then we can take the tensor
product $B \otimes_A C$ of $B$ and $C$ as $A$-modules and this is again an $A$-algebra.
That is, we define the product by
\[
  \Big(\sum b_i \otimes c_i\Big)\Big(\sum b_j' \otimes c_j'\Big) = \sum b_ib_j' \otimes c_ic_j',
\]
and the ring homomorphism $A \to B \otimes_A C$ by $a \mapsto a \otimes 1 (= 1 \otimes a)$. The
fact that the above product is well-defined can easily be seen using the
bilinearity of $bb' \otimes cc'$ with respect to both $(b, c)$ and $(b', c')$. The algebra
$B \otimes_A C$ contains $B \otimes 1$ (short for the subset $\{b \otimes 1 \mid b \in B\} \subset B \otimes C$)
and $1 \otimes C$ as subalgebras, and is generated by these. Note that $B \otimes 1$ is not
necessarily isomorphic to $B$.

\begin{example*}
If $\mathfrak{a}$ is an ideal of $A$ and $C = A/\mathfrak{a}$, then
$B \otimes_A C = B \otimes_A (A/\mathfrak{a}) = B/\mathfrak{a}B$, and the above $B \otimes 1$ is also equal to
$B/\mathfrak{a}B$.
\end{example*}

\begin{example*}
If $B$ is an $A$-algebra and $A[X]$ is the polynomial ring over $A$ then
$B \otimes_A A[X]$ can be identified with $B[X]$. Indeed, $A[X]$ is a free $A$-module
with basis $\{X^\nu \mid \nu = 0, 1, 2, \ldots\}$, so that $B \otimes_A A[X]$ is also the free $B$-module
with basis $\{X^\nu\}$, and is isomorphic to $B[X]$ both as an $A$-module and as a
ring. Similarly for the polynomial ring in several variables.
\end{example*}

\medskip
\noindent\textbf{Direct limits}

A \emph{directed set} is a partially ordered set $\Lambda$ such that for any
$\lambda, \mu \in \Lambda$ there exists $\nu \in \Lambda$ with $\lambda \le \nu$ and $\mu \le \nu$. For example, a
totally ordered set is directed; the set of finite subsets of a set $S$, ordered by
inclusion, is a directed set which is not totally ordered.

Suppose that for each element $\lambda$ of a directed set $\Lambda$ we are given a set $M_\lambda$,
and whenever $\lambda \le \mu$ we are given a map $f_{\mu\lambda}\colon M_\lambda \to M_\mu$ satisfying the
conditions
\[
  f_{\lambda\lambda} = 1, \quad \text{and} \quad f_{\nu\mu} \circ f_{\mu\lambda} = f_{\nu\lambda}
      \quad \text{for } \lambda \le \mu \le \nu;
\]
we express all this data as $(M_\lambda; f_{\mu\lambda})$, and refer to it as a \emph{direct system}
over $\Lambda$ (or indexed by $\Lambda$). If each $M_\lambda$ is an $A$-module, and each $f_{\mu\lambda}$
$A$-linear we speak of a direct system of $A$-modules; if each $M_\lambda$ is a ring,
and each $f_{\mu\lambda}$ a ring homomorphism, a direct system of rings. More generally,
we can define direct systems in any category.

Given two direct systems $\mathcal{F} = \{M_\lambda; f_{\mu\lambda}\}$ and
$\mathcal{F}' = \{M'_\lambda; f'_{\mu\lambda}\}$ indexed by the same set, a morphism
$\varphi\colon \mathcal{F} \to \mathcal{F}'$ is a system of maps $\{\varphi_\lambda\colon M_\lambda \to M'_\lambda\}$ such that
\[
  f'_{\mu\lambda} \circ \varphi_\lambda = \varphi_\mu \circ f_{\mu\lambda} \quad \text{for } \lambda \le \mu.
\]
By a map from $\mathcal{F}$ to a set $X$ we mean a system $\{\varphi_\lambda\}$ of maps
$\varphi_\lambda\colon M_\lambda \to X$ satisfying $\varphi_\lambda = \varphi_\mu \circ f_{\mu\lambda}$ for $\lambda \le \mu$. Now if a
map $\psi\colon \mathcal{F} \to M_\infty$ from $\mathcal{F}$ to a set $M_\infty$ has the universal property for maps
from $\mathcal{F}$ to sets, that is, if for any map $\varphi\colon \mathcal{F} \to X$ there exists a unique
map $h\colon M_\infty \to X$ such that $\varphi_\lambda = h \circ \psi_\lambda$ for all $\lambda \in \Lambda$, then $M_\infty$ is called
the \emph{direct limit} of $\mathcal{F}$, or simply the limit of $\mathcal{F}$, and we write
$M_\infty = \varinjlim M_\lambda$, or $M_\infty = \lim M_\lambda$. As one sees easily from the definition, a map
$\varphi\colon \mathcal{F} \to \mathcal{F}'$ induces a map $\lim M_\lambda \to \lim M'_\lambda$, which in this book we write
$\varphi_\infty$ or $\lim \varphi$.

The limit of a direct system $\mathcal{F} = \{M_\lambda; f_{\mu\lambda}\}$ always exists. In order to
construct it we do the following: take the disjoint union $\bigsqcup_\lambda M_\lambda$ of the
$M_\lambda$, and define a relation $\sim$ by
\[
  x \sim y \quad (x \in M_\lambda, y \in M_\mu) \iff \text{there exists a } \nu
      \text{ with } \lambda \le \nu, \mu \le \nu \text{ and } f_{\nu\lambda}(x) = f_{\nu\mu}(y).
\]
Then it is easy to see that $\sim$ is an equivalence relation. We write $M_\infty$ for the
quotient set $(\bigsqcup_\lambda M_\lambda)/\sim$, that is the set of equivalence classes under $\sim$;
then one sees easily that $M_\infty$ satisfies the conditions for a direct limit. We write
$\lim x \in M_\infty$ for the equivalence class of $x \in M_\lambda$. If $\mathcal{F}$ is a direct system of modules
then $M_\infty$ can be given a natural structure of $A$-module, and
$x \mapsto \lim x$ is an $A$-linear map from $M_\lambda$ to $M_\infty$. Similarly for direct systems
or rings.

The above is general theory. In this book the following two theorems are of
particular importance.

\begin{theorem}\label{thm:A1}
Let $A$ be a ring, $N$ an $A$-module, and let $\mathcal{F} = \{M_\lambda; f_{\mu\lambda}\}$ be a direct
system of $A$-modules. Then
\[
  \varinjlim (M_\lambda \otimes_A N) = (\varinjlim M_\lambda) \otimes_A N.
\]
(In other words, tensor product commutes with direct limits.)
\end{theorem}

\begin{proof}
Set $\lim M_\lambda = M_\infty$ and $\lim(M_\lambda \otimes N) = L_\infty$. We write
$\varphi_\lambda\colon M_\lambda \to M_\infty$ for the $A$-linear map given by $x \mapsto \lim x$, so that
$\{\varphi_\lambda \otimes 1\}$ is a map from the direct system $(M_\lambda \otimes N; f_{\mu\lambda} \otimes 1)$ to the
$A$-module $M_\infty \otimes N$; this determines a unique $A$-linear map
$h\colon L_\infty \to M_\infty \otimes N$. For $x \in M_\lambda$ and $y \in N$ we have
$h(\lim(x \otimes y)) = (\lim x) \otimes y$. On the other hand, fixing $y \in N$ we can define
$g_{\lambda, y}\colon M_\lambda \to L_\infty$ by $g_{\lambda, y}(x) = \lim(x \otimes y)$, and in the limit we get
\[
  g_y\colon M_\infty \to L_\infty.
\]
If $x_\infty \in M_\infty$ we can write $x_\infty = \lim x$ for some $\lambda$ and some $x \in M_\lambda$. Then
$g_y(x_\infty) = g_{\lambda, y}(x) = \lim(x \otimes y)$. From this we can see that $g_y(x_\infty)$ is
bilinear in $x_\infty$ and in $y$, and so defines an $A$-linear map
$g\colon M_\infty \otimes N \to L_\infty$ such that $g(x_\infty \otimes y) = g_y(x_\infty)$. Now it is easy to see
that $g$ and $h$ are inverse maps, so that $M_\infty \otimes N \cong L_\infty$.
\end{proof}

\begin{theorem}\label{thm:A2}
Suppose that we have three direct systems of $A$-modules indexed by the same
set $\Lambda$, $\mathcal{F}' = \{M'_\lambda; f'_{\mu\lambda}\}$, $\mathcal{F} = \{M_\lambda; f_{\mu\lambda}\}$ and
$\mathcal{F}'' = \{M''_\lambda; f''_{\mu\lambda}\}$, and maps $\{\varphi_\lambda\}\colon \mathcal{F}' \to \mathcal{F}$ and
$\{\psi_\lambda\}\colon \mathcal{F} \to \mathcal{F}''$ such that for every $\lambda$,
\[
  M'_\lambda \xrightarrow{\varphi_\lambda} M_\lambda \xrightarrow{\psi_\lambda} M''_\lambda
\]
is an exact sequence; then the sequence obtained in the limit
\[
  \varinjlim M'_\lambda \xrightarrow{\varphi_\infty} \varinjlim M_\lambda \xrightarrow{\psi_\infty} \varinjlim M''_\lambda
\]
is also exact. (In other words, direct limit is an exact functor.)
\end{theorem}

\begin{proof}
Write $M_\infty$ for the limit of $\mathcal{F}$, and let $y_\infty \in M_\infty$ be such that
$\psi_\infty(y_\infty) = 0$. For some $\lambda$ and $y \in M_\lambda$ we can write $y_\infty = \lim y$, and then
$0 = \psi_\infty(\lim y) = \lim \psi_\lambda(y)$, so that for some $\mu \ge \lambda$ we have
$\psi_\mu(f_{\mu\lambda}(y)) = 0$; the left-hand side here is equal to
$\psi_\mu(f_{\mu\lambda}(y))$, so that by assumption there is $x \in M'_\mu$ such that
$f_{\mu\lambda}(y) = \varphi_\mu(x)$. Thus $y_\infty = \lim f_{\mu\lambda}(y) = \lim \varphi_\mu(x) =
\varphi_\infty(\lim x) \in \Im(\varphi_\infty)$. Also, $\psi_\infty \circ \varphi_\infty = 0$ is obvious.
\end{proof}

Given an $A$-module $M$, write $\{M_\lambda\}_{\lambda \in \Lambda}$ for the collection of all finitely
generated submodules of $M$. We define a partial order on $\Lambda$ by letting
$\lambda \le \mu$ if $M_\lambda \subset M_\mu$; this makes $\Lambda$ into a directed set, and we write
$f_{\mu\lambda}\colon M_\lambda \to M_\mu$ for the natural inclusion. Then $\{M_\lambda; f_{\mu\lambda}\}$ is a
direct system of $A$-modules, the limit of which is the original $M$, that is
$M = \varinjlim M_\lambda$. Hence any $A$-module can be expressed as a direct limit of
finitely generated $A$-modules.

In a similar way, given any ring $A$ and a subring $A_0 \subset A$, we can express
$A$ as the direct limit of subrings which are finitely generated over $A_0$ as
rings. If we take $A_0$ to be the minimal subring of $A$ (that is, the image in $A$
of $\mathbb{Z}$), then a ring which is finitely generated over $A_0$ is Noetherian, and
hence every ring is a direct limit of Noetherian rings.

\medskip
\noindent\textbf{Inverse limits}

Inverse systems and inverse limits are defined as the dual notions to direct
systems and direct limits, that is by reversing all the arrows in the
definitions. That is, we take a directed set $\Lambda$ as indexing set; an inverse
system of sets is the data of a set $M_\lambda$ for each $\lambda \in \Lambda$, and of a map
$f_{\lambda\mu}\colon M_\mu \to M_\lambda$ whenever $\lambda \le \mu$, such that
\[
  f_{\lambda\lambda} = 1, \quad \text{and} \quad f_{\lambda\mu} \circ f_{\mu\nu} = f_{\lambda\nu}
      \quad \text{for } \lambda \le \mu \le \nu;
\]
we write this as $(M_\lambda; f_{\lambda\mu})$. A morphism between two inverse systems with
the same indexing set, and a map from a set $N$ to an inverse system
$\mathcal{F} = \{M_\lambda; f_{\lambda\mu}\}$ are defined dually to the case of direct systems. We say
that $M_\infty$ is an \emph{inverse limit}, or projective limit of $\mathcal{F}$, and write
$M_\infty = \varprojlim M_\lambda$, if there is a map $\varphi = \{\varphi_\lambda\}\colon M_\infty \to \mathcal{F}$ which has the
property that for any set $X$, and any map $\psi = \{\psi_\lambda\}\colon X \to \mathcal{F}$, there exists a
unique map $h\colon X \to M_\infty$ such that $\psi_\lambda = \varphi_\lambda h$ for all $\lambda$. To prove the
existence of $\varprojlim M_\lambda$ we only have to let $M_\infty$ be the following subset of the
direct product $\prod_\lambda M_\lambda$:
\[
  M_\infty = \{(x_\lambda)_{\lambda \in \Lambda} \mid \lambda \le \mu \Rightarrow x_\lambda =
      f_{\lambda\mu}(x_\mu)\}.
\]
If each $M_\lambda$ is a module and each $\varphi_{\lambda\mu}$ is a linear map then this $M_\infty$ is a
submodule of the direct product module, and is the inverse limit of modules.
In a similar way, the inverse limit of an inverse system of rings is again a
ring.

\begin{example*}
Let $\Lambda = \{1, 2, 3, \ldots\}$ and let $p$ be a prime number. Consider the inverse
system of rings
\[
  \mathbb{Z}/(p) \leftarrow \mathbb{Z}/(p^2) \leftarrow \mathbb{Z}/(p^3) \leftarrow \cdots,
\]
where each arrow is the natural homomorphism. The inverse limit
$\varprojlim \mathbb{Z}/(p^n)$ is known as the ring of \emph{$p$-adic integers}. Its elements are of
the form
\[
  (a_1, a_2, a_3, \ldots), \quad \text{with } a_i \in \mathbb{Z}/(p^i) \text{ and } a_i = a_{i-1} \bmod p^{i-1};
\]
addition and multiplication is carried out term-by-term:
\[
  (a_1, a_2, \ldots) + (b_1, b_2, \ldots) = (a_1 + b_1, a_2 + b_2, \ldots)
\]
\[
  (a_1, a_2, \ldots) \cdot (b_1, b_2, \ldots) = (a_1b_1, a_2b_2, \ldots).
\]
More generally, if $A$ is any ring and $I$ an ideal of $A$, the inverse limit
$\varprojlim A/I^n$ of the inverse system of rings $A/I \leftarrow A/I^2 \leftarrow \cdots$ is called the
\emph{$I$-adic completion} of $A$ (see \S\ref{sec:8}).
\end{example*}

Taking the inverse limit of an inverse system of modules is a left-exact
functor, but is not an exact functor, so that the analog of
Theorem~\hyperref[thm:A2]{A2} for inverse systems does not hold.

\begin{example*}
Consider the diagram
\[
\begin{tikzcd}
0 \arrow[r] & \mathbb{Z} \arrow[r, "n"] \arrow[d] & \mathbb{Z} \arrow[r] \arrow[d] & \mathbb{Z}/(n) \arrow[r] \arrow[d] & 0 \\
0 \arrow[r] & \mathbb{Z} \arrow[r, "p"] \arrow[d] & \mathbb{Z} \arrow[r] \arrow[d] & \mathbb{Z}/(n) \arrow[r] \arrow[d] & 0 \\
0 \arrow[r] & \mathbb{Z} \arrow[r, "n"] & \mathbb{Z} \arrow[r] & \mathbb{Z}/(n) \arrow[r] & 0
\end{tikzcd}
\]
here $p$ and $n$ are coprime integers, and the arrows marked $\times$ are multiplication
by $n$. The rows are exact sequences, and each column defines an inverse
system. The left-hand and middle columns have $0$ as their inverse limits, but
since every arrow in the right-hand column is an isomorphism, the inverse
limit is isomorphic to $\mathbb{Z}/n$. Thus going to the inverse limit, we find that
$0 \to 0 \to \mathbb{Z}/(n)$ is exact, but the second arrow is not surjective.
\end{example*}

% ============ 习题 ============
\section*{Exercises to Appendix A}
\emph{Prove the following propositions.}

\noindent\textbf{A.1.}\quad Let $M$ and $N$ be $A$-modules. If the natural map
$M' \otimes_A N \to M \otimes_A N$ is injective for every finitely generated submodule
$M' \subset M$ then the same thing holds for every submodule $M' \subset M$.

\medskip
\noindent\textbf{A.2.}\quad Let $A$ be a ring, and $B$, $C$, $D$ (commutative) $A$-algebras. Then to
give a homomorphism of $A$-algebras from $B \otimes_A C$ to $D$ is the same thing as to
give a pair of homomorphisms of $A$-algebras $B \to D$ and $C \to D$; in other words,
$B \otimes_A C$ is the category-theoretical direct product of $B$ and $C$ in the category of
$A$-algebras.
